\subsection{Analisis Kebutuhan Offline First dan Batching Mechanism}

Ini menjawab RM2 di level analisis. Gunakan literatur offline-first:

\begin{enumerate}
    \item Paper / artikel arsitektur offline-first untuk sistem enterprise \& healthcare:
    \begin{itemize}
        \item full data sync vs action queue.

        \item case study field apps healthcare/banking.
    \end{itemize}

    \item Kaitkan dengan PGHD:
    \begin{itemize}
        \item pasien mungkin di area sinyal buruk; data harus tetap terkumpul.
        
        \item offline-first bukan fitur tambahan, tapi requirement.
    \end{itemize}
\end{enumerate}

Hal yang bisa kamu diskusikan:
\begin{enumerate}
    \item Dua pendekatan offline-first: full sync (local mirror server state) vs action queue (antrian data/write).
    
    \item Untuk kasusmu (PGHD → mostly write-heavy: log data sensor): action queue / event log (data queuing dan batching) lebih cocok, tidak perlu full mirror state dari server.
\end{enumerate}

Dari sini kamu turunkan kebutuhan:
\begin{enumerate}
    \item O1: harus ada local storage (DB) untuk log PGHD,
    
    \item O2: harus ada mekanisme penjadwalan pengiriman (batch) ketika koneksi tersedia,
    
    \item O3: harus ada mekanisme penanganan kegagalan (retry, penanda status).
\end{enumerate}